\chapter{Cryptographic Protocols}
This lecture is about the security of message protocols.
There are different definitions of security, in the course of this lecture we define, what security means in this context of cryptographic protocols, and how to automatically verify the security of a protocol.

\section{Foundations}
First and foremost, we define the cryptographic primitives and basic operations.

\newcommand{\encs}{\texttt{enc}^\text{s}}
\newcommand{\enca}{\texttt{enc}^\text{a}}
\newcommand{\decs}{\texttt{dec}^\text{s}}
\newcommand{\deca}{\texttt{dec}^\text{a}}
\begin{definition}[Cryptographic Primitives: Encryption]
    For \(X,Y,K \subseteq \Sigma^*\) for alphabet \(\Sigma\), we define the \emph{symmetric} encryption \(\encs : X \times K \rightarrow Y\) and decryption \(\decs : Y \times K \rightarrow X\) as some bijective functions with
    \[\decs_k \circ \encs_k = x \mapsto x.\]
    Furthermore, for \(\hat K \subseteq \Sigma^*\) and some bijective mapping to \(K\), we define the \emph{asymmetric} encryption \(\enca : X \times K \rightarrow Y\) and decryption \(\deca : Y \times \hat K \rightarrow X\) functions with property
    \[\deca_{\hat{k}} \circ \enca_{k} = x \mapsto x.\]
    In both encryption primitives, we assume the security property that there is no easily computable direct mapping \(Y \rightarrow X\), nor \(K \rightarrow \hat K\).

\end{definition}

\newcommand{\mac}{\texttt{mac}}
\newcommand{\tests}{\texttt{test}^\text{s}}
\newcommand{\testa}{\texttt{test}^\text{a}}
\newcommand{\ok}{\texttt{ok}}
\newcommand{\err}{\texttt{error}}
\newcommand{\sig}{\texttt{sig}}
\begin{definition}[Cryptographic Primitives: Signatures]
    For \(X,K,T \subseteq \Sigma^*\), we define the symmetric signature \(\mac : X \times K \rightarrow T\) and verification \(\tests : X \times K \times T \rightarrow \{\ok, \err\}\), such that
    \[\tests_k(x, \mac_k(x)) \overset{!}{=} \ok.\]
    Furthermore, for \(\hat K \subseteq \Sigma^*\) and some bijective mapping to \(K\), we define the asymmetric signature \(\sig : X \times \hat K \rightarrow T\) and verification \(\testa : X \times K \times S \rightarrow \{\ok, \err\}\), such that
    \[\testa_k(x, \sig_{\hat{k}}(x)) \overset{!}{=} \ok.\]
    Again, we assume the security property, that there is no easily computable direct mapping \(X \rightarrow T\), nor \(K \rightarrow \hat K\).

    For signatures, we often write \(\sig_k(t)\) instead of \(\sig_{\hat k}(t)\) to highlight, that the public key \(k\) is known by all, even though the private key is needed to sign \(t\).
\end{definition}

\newcommand{\hash}{\texttt{hash}}
\begin{definition}[Cryptographic Primitives: Hashing]
    For \(X, T \subseteq \Sigma^*\), we define the bijective mapping \(\hash : X \rightarrow T\), where \(\hash^{-1}\) is not easily computable.
\end{definition}

\begin{definition}[Cryptographic Primitives: Nonces and Random Numbers]
    Lastly, we assume, that we can generate perfectly random numbers or messages we call nonces, typically denoted by \(N\).
\end{definition}

\section{Formal Protocol Model}
Because we have difficulties when talking about security in an ambiguous model, we have to define a formal model to represent protocols, before defining what security even means.

In the context of this lecture, a formal model must support the following features:
1. The adversary must be able to control the routing of messages.
2. The model must be able to generate nonces.
3. The adversary may send arbitrary messages, limited only by cryptography.
4. The adversary must be able to control the scheduling.

Firstly, we have to define, how to notate messages and terms.

\newcommand{\T}{\mathcal{T}}
\newcommand{\C}{\mathcal{C}}
\newcommand{\V}{\mathcal{V}}
\newcommand{\ID}{\text{IDs}}
\begin{definition}[Terms and Messages]\label{def:terms}
    We denote the set of terms \(\T\) as the smallest set with the following properties for \(\C\) the set of constants, \(\V\) the set of variables, and \(\ID\) the set of names.
    \begin{IEEEeqnarray*}{r'c'c'c'l}
        & & \{\epsilon\} \cup \C \cup \V \cup \ID & \subseteq & \T \\
        i \in \N, a \in \ID & \Rightarrow & N_i, N_i^a, k_a, \hat k_a & \in & \T \\
        t_1, t_2 \in \T & \Rightarrow & [t_1, t_2] & \in & \T \\
        t, k \in \T & \Rightarrow & \encs_k(t) & \in & \T \\
        t \in \T \wedge a \in \ID & \Rightarrow & \enca_{k_a}(t) & \in & \T \\
        t, k \in \T & \Rightarrow & \mac_k(t) & \in & \T \\
        t \in \T \wedge a \in \ID & \Rightarrow & \sig_{k_a}(t) & \in & \T \\
        t \in \T & \Rightarrow & \hash(t) & \in & \T
    \end{IEEEeqnarray*}
    We relax this definition by allowing \(N_A\) or \(N_A''\) to be nonces generated by \(A \in \ID\) or writing \([t_0, [t_1, t_2]]\) as \(t_0, t_1, t_2\).
    Terms can be represented as a tree or a directed acyclic graph (DAG).

    We call terms without variables a \emph{ground term} or a \emph{message}.
\end{definition}

In the tree or DAG of a term, we can trace to some position by \emph{undoing} operators specified in \cref{def:terms}, we call these terms \emph{subterms}.

\newcommand{\Sub}{\text{Sub}}
\begin{definition}[Subterms]
    For some term \(t \in \T\) and some position \(p \in \N^+\), we write \(t|_p\) as the subterm reached in \(t\) by following \(p\).

    More formally, we define the set of subterms \(\Sub(t)\) of term \(t \in \T\), as the smallest set with the following properties.
    \begin{IEEEeqnarray*}{r'c'c'c'l}
        & & t & \in & \Sub(t) \\
        t' \in \Sub(t) & \Rightarrow & \Sub(t') & \subseteq & \Sub(t) \\
        \exists t_1, t_2 : t = [t_1, t_2] & \Rightarrow & t_1, t_2 & \in & \Sub(t) \\
        \exists k, t' : t = \encs_k(t') & \Rightarrow & k, t' & \in & \Sub(t) \\
        \exists k, t' : t = \enca_k(t') & \Rightarrow & k, t' & \in & \Sub(t) \\
        \exists k, t' : t = \sig_k(t') & \Rightarrow & k, t' & \in & \Sub(t) \\
        \exists k, t' : t = \mac_k(t') & \Rightarrow & k, t' & \in & \Sub(t) \\
        \exists t' : t = \hash(t') & \Rightarrow & t' & \in & \Sub(t)
    \end{IEEEeqnarray*}
    We also define the set of subterms of a set of terms \(S \subseteq \T\) as
    \[\Sub(S) := \bigcup_{t \in S} \Sub(t).\]
    Finally, we denote \(\Sub^1(S)\) as the set of direct subterms of \(t\).
\end{definition}

As mentioned, the adversary may generate messages only limited by cryptography.
Therefore they should not be able to send a message containing a non-derivable term.
Intuitively, we define the \emph{derivability} of some message as the \emph{Dolev-Yao closure}.

\newcommand{\DY}{\text{DY}}
\begin{definition}[Dolev-Yao Closure]
    For some \(S \subseteq \T\), we define the Dolev-Yao closure of \(S\), denoted \(\DY(S)\) as the set of messages, that are derivable from \(S\).
    The Dolev-Yao Closure of \(S\) is the smallest set with the following properties.
    \begin{IEEEeqnarray*}{r'c'r}
        S \cup \{\epsilon\} \cup \ID & \subseteq & \DY(S) \\
        {[t_1, t_2]} \in \DY(S) & \Leftrightarrow & t_1, t_2 \in \DY(S) \\
        a \in \ID, t \in \DY(S) & \Rightarrow & \enca_{k_a}(t) \in \DY(S) \\
        t, k \in \DY(S) & \Rightarrow & \encs_k(t), \mac_k(t) \in \DY(S) \\
        t, \hat k \in \DY(S) & \Rightarrow & \sig_k(t) \in \DY(S) \\
        \encs_k(t), k \in \DY(S) & \Rightarrow & t \in \DY(S) \\
        \enca_{k_a}(t), \hat k_a \in \DY(S) & \Rightarrow & t \in \DY(S) \\
        \sig_k(t) \in \DY(S) & \Rightarrow & t \in \DY(S) \\
        \mac_k(t) \in \DY(S) & \Rightarrow & t \in \DY(S) \\
        t \in \DY(S) & \Rightarrow & \hash(t)
    \end{IEEEeqnarray*}
    It should be highlighted, that this definition of the Dolev-Yao closure allows composited keys, e.g., \(\encs_{N_A,N_B}(\texttt{"Alice loves Bob"})\) to communicate a secret between Alice and Bob.
    In this example, the message is derivable from \(S \subseteq \T\), iff \(N_A, N_B \in \DY(S)\), which is the case if both nonces are derivable or if \(N_A, N_B \in S\).
\end{definition}

Now, as seen pretty quickly, the Dolev-Yao closure is always infinite, therefore we can not compute it, but we can go around this problem by handling it as a decision problem.

\newcommand{\Derive}{\textsf{DERIVE}}

\begin{problem}[Derivation of a Term]\label{prob:derive}
    The \(\Derive\) problem has the input \(S \subseteq \T, m \in \T\) and answers if \(m \in \DY(S)\) or not.
    \iffalse
    \begin{IEEEeqnarray*}{l'l}
        \textit{Problem}    & \text{DERIVE} \\
        \textit{Input}      & S \subseteq \T, m \in \T \\
        \textit{Question}   & m \in \DY(S)
    \end{IEEEeqnarray*}
    \fi
\end{problem}

It turns out, \(\Derive\) is decidable in polynomial time using a fixed-point algorithm.
The algorithm finds a shortest derivation of \(t\) iff \(t \in \DY(S)\).
This restricts the search space by using \emph{composition} and \emph{decomposition} rules which allow to derive to and from some term.

\begin{definition}[Composition and Decomposition Rules]
    For some term \(m \in \T\), we define composition and decomposition rules as functions, that describe how \(m\) can be decomposed or composed---what to derive to with \(m\) and how to derive to \(m\).

    We denote decomposition and composition rules as \(L_d(m), L_c(m) : R \rightarrow O\), where \(R, O \subseteq \T\) are the required and obtained terms.
    This is also often written without set brackets, e.g., \(L_c([a,b]) : a,b \rightarrow [a,b]\).

    A whole derivation \(D(m)\) of a message \(m\) from \(S \subseteq \T\) is written as a chain of derivation rules, e.g.,
    \[S = S_0 \rightarrow_{L_0} \rightarrow S_1 \rightarrow_{L_1} \dots \rightarrow S_{n-1} \rightarrow_{L_{n-1}} \rightarrow S_n \ni m.\]
\end{definition}

The success in this notation lies in the fact, that any required non-atomic terms are always direct subterms of the term the rule is applied on/to.
More formally,
\begin{IEEEeqnarray*}{r'c'l}
    L_x(t) : R \rightarrow O & \Rightarrow & R \setminus \T_{\text{at}} \subseteq \{t\} \cup \Sub^1(t)
\end{IEEEeqnarray*}
Composition rules have exactly the terms of \(\Sub^1(m)\) as required terms.
Decomposition rules obtain only terms from \(\Sub^1(m)\).

To reduce the search space even further, we can look at the properties of some shortest derivation \(D_S(m)\) of \(m \in \T\).
\begin{IEEEeqnarray*}{r'c'l}
    L_d(t) \in D_S(m) & \Rightarrow & t \in \Sub(S) \\
    L_c(t) \in D_S(m) & \Rightarrow & t \in \Sub(S \cup \{m\})
\end{IEEEeqnarray*}
To derive \(m\) from \(S\), we only need decompositions of subterms from \(S\) or compositions of subterms of \(S\) or subterms of \(m\).

From these rules, we find the following fixed-point algorithm to solve \(\Derive\) (\cref{prob:derive}) in polynomial time.

\pagebreak % FIXME: Remove when some text is added to the current page.

\begin{algorithm}[Algorithm for \(\Derive\)] \hfill
\begin{lstlisting}
def derive($S$, $m$):
    $S_{\text{old}} = \emptyset$
    while $S_{\text{old}} \ne S$ do
        $S_{\text{old}} = S$
        if $\exists t \in \Sub(S \cup \{m\}) \setminus S, S \rightarrow S \cup \{t\}$ then
            $S = S \cup \{t\}$
        fi
    od
    if $m \in S$ then
        accept
    else
        reject
    fi
\end{lstlisting}
\end{algorithm}

\begin{definition}[Receive/Send Action]
    We define a receive/send (r/s) action as some pair \((r, s) \in \T^2\) which we write \(r \rightarrow s\).
\end{definition}

Currently, we still have some problems, because a protocol should be independent of its participants and needs features not yet implemented.
For example, a protocol with some action \(\enca_{k_B}(A, x) \rightarrow \enca_{k_A}(x, N_B)\) assumes, that we have some way to handle variables, in this case \(x\), and hard codes some recognition of key and nonce origin.

To fix these issues, we introduce \emph{substitutions}.

\begin{definition}[Substitutions]
    Substitutions are functions \(\sigma : \V \rightarrow \T\) with \(\exists^{<\infty} x : \sigma(x) \ne x\).
    A substitution \(\sigma\) is called \emph{ground substitution}, if \(\forall x \in \V, \sigma(x) \ne x: \sigma(x) \text{ is message}\).

    We interpret a situation \(\sigma(x) = x\) as the case, that \(x\) is uninitialized.
    Else, we evaluate inductively for arbitrary terms\footnote{\(\T_{\text{nat}}\) is supposed to be \(\T \setminus \T_{\text{at}}\), aka., non-atomic terms.}.
    \begin{IEEEeqnarray*}{l"r'c'l}
        \text{IB: } x \in \V              & \sigma(x) & & \text{ is already defined} \\
        \text{IS: } x \in \T_{\text{nat}} \setminus \V & \sigma(f(t_1, \dots, t_n)) & := & f(\sigma(t_1), \dots, \sigma(t_n))
    \end{IEEEeqnarray*}
\end{definition}

Putting this together, we can model a protocol step, as a substitution state \(\sigma\), next r/s action, incoming message \(m\), updated substitution \(\sigma \mapsto \sigma'\), and reply message \(m'\).

\begin{definition}[Matching]
    We say a term \(t \in \T\) matches with message \(m\) and substitution \(\sigma\) via substitution \(\sigma'\), if
    \begin{IEEEeqnarray*}{r'c'l}
        \sigma'(t) & = & m \\
        \forall x, \sigma(x) \ne x : \sigma'(x) & = & \sigma(x).
    \end{IEEEeqnarray*}
\end{definition}

\subsection{Formal Protocol Model}
Now we have all important steps to compose a formally sound protocol model.

\newcommand{\I}{\mathcal I}

\begin{definition}[Protocol Instance]
    A protocol instance is a \(n \in \N\) long sequence \(\I \in (\T^2)^n\) of r/s actions
    \[\I = r_0 \rightarrow s_0, r_1 \rightarrow s_1, \dots, r_{n-1} \rightarrow s_{n-1}\]
    with \(\forall i \in n : \V(s_i) \subseteq \bigcup_{j \le i} \V(r_j)\), which basically means, that all variables of some response are already defined in previous received terms.
\end{definition}

\begin{definition}[Protocol]
    A protocol consists of a finite number of instances \(\I_0, \dots, \I_{m-1}\) and a finite set \(I \subseteq \T\) of messages, which models the initial adversary knowledge.
\end{definition}

\subsubsection{Needham-Schroeder Protocol}\label{sub:needham-schroeder}
Using these definitions we can now finally define the Needham-Schroeder protocol using our formal notation.
\begin{IEEEeqnarray*}{r'c'l}
    \I_0 & & \text{Client (Alice)} \\
    \epsilon & \rightarrow & \enca_{k_B}(A, N_A) \\
    \enca_{k_A}(N_A, y) & \rightarrow & \enca_{k_B}(y) \\ \\
    I_1 & & \text{Server (Bob)} \\
    \enca_{k_B}(A, x) & \rightarrow & \enca_{k_A}(x, N_B) \\ \\
    I & & \text{Initial adversary knowledge} \\
    I & = & \{A, B, C, k_A, k_B, k_C, \hat k_C, N_C^1, N_C^2, N_C^3\}
\end{IEEEeqnarray*}
Now we need to model, how the adversary can attack a protocol.
In this case, the protocol is insecure, because, if \(A\) starts the protocol with \(C\), \(C\) can redirect the message to \(B\), who thinks that \(A\) wants to talk with \(B\) and sends the response to \(A\).
Finally, \(A\) sends \(N_B\) to \(C\), thinking that \(N_B\) was generated by \(C\) and \(C\) redirects \(N_B\) again to \(B\).
Now the secret \(N_A \oplus N_B\) is known by \(A\), \(B\), and \(C\) even though \(B\) thinks, that only \(A\) also knows \(N_A\) and \(N_B\).
Thus, \(B\) may share secrets using the composited key \(N_A \oplus N_B\) unbeknownst to \(C\) and not \(A\).

\subsection{Modeling Protocol Attacks}
The current notation hard codes, that \(A\) and \(B\) want to talk with each other.
Currently, the case that \(A\) wants to talk with \(C\) and \(B\) with \(A\) is not possible to model.

Additionally, the adversary must be able to control the scheduling, which is not yet implemented.

\begin{definition}[Execution Order]
    For protocol \(P := (\{\I_0, \dots, \I_{n-1}\}, I)\), we say each instance \(I_\cdot\) has \(|I_\cdot|\) steps.
    An \emph{execution order} for \(P\) is a sequence\footnote{In this lecture we define \([n]\) as \(\{0, \dots, n-1\}\)} \(o \in [n]^+\) such that each \(j \in [n]\) appears at most \(|I_j|\) times.
    We denotate \(o(t)\) as the \(t\)-th element in \(o\) and \(\#o(t) := |\{l \mid l < t \wedge o(l) = o (t)\}|\).

    \(o(t)\) can be seen as the current instance that gets \emph{handled}, which makes \(o\) itself a sequence of instance steps.
    Therefore for some timestamp \(t\), \(\#o(t)\) is the current state of the referenced instance.
\end{definition}

\begin{definition}[Protocol Run / Attack]
    For protocol \(P := (\{\I_0, \dots, \I_{n-1}\}, I)\) and all \(i\)-th rules of the \(j\)-th instance \(\I_j\) being named\footnote{We will reuse this notation.} \(r_i^j \rightarrow s_i^j\), we define a \emph{protocol run} of \(P\) or \emph{attack} as an execution order \(o\), and a substitution \(\sigma\), such that
    \[\forall t \in [|o|] : \sigma(r_{\#o(t)}^{o(t)}) \in \DY\left(I \cup \left\{\sigma(s_{\#o(l)}^{o(l)}) \mid l \in [t]\right\}\right).\]
    In layman's terms, the substitution is defined such, that the values it can know from the received terms are derivable from previous send terms and the initial adversary knowledge.
    This is really practical, as the adversary knows every message that is send and received.
\end{definition}

Now we still need to define, what a \emph{successful} attack is, which is dependant on the definition of security, which is appropriate to the relevant context.
As we discuss the security of protocols, we use the following definition.

\newcommand{\fail}{\texttt{FAIL}}
\newcommand{\brea}{\texttt{BREAK}}

\begin{definition}[Successful Attack]
    We call an attack on some protocol \(P\) successful, if
    \[\fail \in \DY\left(I \cup \left\{\sigma(s_{\#o(l)}^{o(l)}) \mid l \in [|o|]\right\}\right).\]
\end{definition}

Now we can just introduce some r/s actions which receives a pair of the \(\brea\) symbol with some hopefully underivable message, which sends \(\fail\) on response.

We can thus extend the protocol definition in \autoref{sub:needham-schroeder} with the following r/s action for Bob, to model the security.
\begin{IEEEeqnarray*}{r'c'l}
    {[\brea, N_B]} & \rightarrow & \fail
\end{IEEEeqnarray*}

Currently, the model is still incomplete, as we still need to implement some things.
For example, we want to keep partners flexible, e.g., allowing \(A \rightarrow C\) and \(A \rightarrow B\) in the same definition.
We also want to allow multiple sessions of the protocol to be able to be simulated at the same time, which currently is not possible.
We will fix both these problems on the way.

\section{Rusinowitch-Turuani Theorem}
Our end goal is the automatic analysis of protocols.
By this definition, this means solving \cref{prob:insecure}.

\newcommand{\Insecure}{\textsf{INSECURE}}
\begin{problem}[Security of a Protocol]\label{prob:insecure}
    The \(\Insecure\) problem has input protocol \(P\) and decides if there is a successful attack on \(P\) or not.
\end{problem}

\newcommand{\NP}{\textsf{NP}}

The way \(\Insecure\) is defined, we quickly run into undecidability for more complicated contexts, but the problem is in fact decidable and \(\NP\)-complete for finite sessions.
The proof of this, the Rusinowitch-Turuani Theorem, is the main goal of this section.

\begin{theorem}[Rusinowitch-Turuani]\label{the:rutu}
    \(\Insecure\) is \(\NP\)-complete.
\end{theorem}

\subsection{Using Directed Acyclic Graphs to Model Terms}

As some attacks may be exponential in length, we have to shorten the representation of terms and attacks to consider \(\NP\) completeness.
Therefore we redefine terms as DAGs.

\newcommand{\DAG}{\textsf{DAG}}

\begin{definition}[Terms as DAGs]
    For some \(S \subseteq \T\), we define the \emph{DAG representation} of \(S\) as an edge-labeled graph \(G := (V, E)\) with
    \begin{IEEEeqnarray*}{r'c'l}
        V & := & \Sub(S) \\
        E & := & \left\{v_s \overset{\text{left}}{\rightarrow} v_e \in V^2 \mid \exists b \in \T: v_s \in \left\{[v_e, b], \encs_{v_e}(b), \enca_{v_e}(b)\right\}\right\} \\
        & \cup & \left\{v_s \overset{\text{right}}{\rightarrow} v_e \in V^2 \mid \exists b \in \T: v_s \in \left\{[b, v_e], \encs_{b}(v_e), \enca_{b}(v_e)\right\}\right\}
    \end{IEEEeqnarray*}
    Moreover, we define \(|S|_\DAG := |V|\).
\end{definition}

The idea of this definition is to show, that if there exists some attack on a protocol \(P\), then there exists an attack \((o, \sigma)\), such that
\[\exists n \in \N : P \mapsto |\{\sigma(x) \mid x \in \V(P)\}|_\DAG \in \mathcal O(P \mapsto |P|^n).\]
We need this, in order to prove \cref{the:rutu}, as the solution substitution \(\sigma\) must not be larger than a polynomial dependant on \(|P|\).
Else, a solution can not be verified in polynomial time, making the problem too hard for \(\NP\)-completeness.

\begin{definition}[Notations for DAG Representations of Terms]
    For \(S \subseteq \T\) and substitution \(\sigma\), we write \(S\sigma := \{\sigma(t) \mid t \in S\}\).
    Additionally, for \(t \in \T, x \in \V\), we write \([x \leftarrow t]\) as substitution with \(x \mapsto t\) and \(y \ne x \Rightarrow y \mapsto y\).
\end{definition}

We need these notations, as thus we see, that for \(S \subseteq \T, x \in \V, t\) message, we have
\[|S[x \leftarrow t]|_\DAG \le |S \cup \{t\}|_\DAG.\]
In other words, applying \(x \leftarrow t\) on an arbitrary amount of variables in \(S\) is not more expensive than adding \(t\) directly to \(S\).

A corollary of this is, that for \(S \subseteq \T\) and \(\sigma\) a ground substitution on \(x_1, \dots, x_k \in \V\), we have
\[|S\sigma|_\DAG \le |S \cup \{\sigma(x_1), \dots, \sigma(x_k)\}|_\DAG.\]
Simply, we can substitute to compress the DAG of \(S\) further.

\subsection{Short Attacks}
As mentioned, we want to show, that for each successful attack, there is a short representation of it, for that we first have to define, what we even consider the size of an attack.

\begin{definition}[Size of an Attack]
    For an attack \((o, \sigma)\) on a protocol \(P\), we define the size \(|\sigma|\) of the attack as \(\sum_{x \in \V(P)}|\{\sigma(x)\}|_\DAG\).

    A minimal attack is subsequently defined as an attack, with no smaller attack possible.
\end{definition}

In most cases we have\footnote{We write \(|\sigma(\cdot)|_\DAG\) for \(|\{\sigma(\cdot)\}|_\DAG\).} \(|\sigma(\cdot)|_\DAG = 1\), as often \(\sigma(\cdot)\) is atomic and thus has no non-trivial subterms.

By the parsing lemma, if \((o, \sigma)\) is a minimal successful attack on \(P\), and there is some variable \(x\) with \(|\sigma(x)|_\DAG > 1\), there is some term \(t\) with\footnote{\(\T(P)\) is the set containing all terms of each r/s action in \(P\).}
\begin{IEEEeqnarray*}{l}
    t \in \Sub(\T(P)) \\
    t \not\in \V \\
    \sigma(t) = \sigma(x)
\end{IEEEeqnarray*}
In other words, complex terms mapped from variables always match with terms in the protocol.
Each term the attacker assigns to a variable is already present in the protocol.

The proof of the parsing lemma is available on \url{https://cloud.rz.uni-kiel.de/index.php/s/3osFeYcC6NwRMij?dir=/&editing=false&openfile=true}.

Finally, we can prove, that we can represent an attack \((o, \sigma)\) on \(P\) in polynomial space, as we can show, that
\[\forall x \in \V : |\sigma(x)|_\DAG \le |\T(P)|_\DAG.\]
Therefore, to prove \cref{the:rutu}, we only need to show, that \(\Insecure\) is \(\NP\)-hard.
Thus, we need to prove, that there is a polynomial reduction from another \(\NP\)-hard problem to \(\Insecure\).

\newcommand{\TSat}{\textsf{3-SAT}}

The lecture shows this by reducing \(\TSat \le_p \Insecure\).
\(\TSat\) is defined\footnote{Hint: I redefined the problem in a way that feels more familiar to the notation used in the lecture.} in \cref{prob:3Sat}.

\newcommand{\true}{\texttt{true}}
\newcommand{\false}{\texttt{false}}

\begin{problem}[\TSat]\label{prob:3Sat}
    The \(\TSat\) problem has input formula \(\phi := \bigwedge_{i \in [n]}(l_1^i \vee l_2^i \vee l_3^i)\) for literals\footnote{\(r_{i,j}\) is the index of \(x\) that \(l_j^i\) is based on.} \(l_j^i \in \{x_{r_{i,j}}, \overline x_{r_{i,j}}\}\) for \(x_{r_{i,j}} \in \V\).
    The problem decides if \(\phi\) is satisfiable, i.e., there exists a substitution \(\sigma : \V \rightarrow \{\true, \false\}\), such that \(\phi \sigma \rightarrow \true\).
\end{problem}

The reduction works by constructing a protocol \(P\) from formula \(\phi\) in the following way.
We create an instance \(A\), which receives an assignment \(I\) from the adversary and distributes across all \(B_j\) instances using symmetric encryption.
Each \(B_j\) now corresponds to \(l_j^1 \vee l_j^2 \vee l_j^3\) and sends \(k_i\) to \(C\), if the literal is satisfiable using \(I\).
There are 7 instances in each \(B_j\), as for three literals, there are 7 possible instances, which satisfy the sentence.
Finally, \(C\) sends \(\fail\), if it receives \(k_1, \dots, k_n\).
The initial adversary knowledge is \(\{\true, \false\}\).

Computing \(P\) from \(\phi\) is obviously possible in polynomial time, and correctly accepts and rejects \(\phi\).
Therefore, \(\Insecure\) is \(\NP\)-hard and thus \(\NP\)-complete.

\section{Automatic Analysis}
As mentioned, the goal of this chapter is the automatic analysis of security protocols.
We quickly run into undecidability if we try to expand our formal model.
E.g., allowing multiple and parallel sessions, makes \(\Insecure\) undecidable.

We also still have to handle, how the adversary controls the routing of messages and who chooses to initialize with whom.

We define a version of \(\Insecure\), which allows parallel sessions in \cref{prob:uninsecure}.

\newcommand{\UnInsecure}{\textsf{UNBOUNDED-INSECURE}}

\begin{problem}[Security of an Unbound Protocol]\label{prob:uninsecure}
    The \(\UnInsecure\) problem has input protocol \(P\) and decides, if there is an insecure protocol based\footnote{In a way such that we only allow extensions of the protocol by copying instances and changing identities.} on \(P\).
\end{problem}

\begin{theorem}[Undecidability of Unrestricted Protocols]\label{the:uninsecure}
    \(\UnInsecure\) is undecidable.
\end{theorem}

\newcommand{\PCP}{\textsf{PCP}}

We show \cref{the:uninsecure} by reducing \(\PCP \le \UnInsecure\).

\newcommand{\0}{\texttt{0}}
\newcommand{\1}{\texttt{1}}

\begin{problem}[Post Correspondence Problem]\label{prob:pcp}
    The \(\PCP\) problem has input two lists of words \(\alpha, \beta \in (\{\0, \1\}^+)^n\), and decides if there is a sequence \(I \in [n]^+\), such that \(\bigtimes_{i \in I}\alpha_i = \bigtimes_{i \in I}\beta_i\).
\end{problem}

\cref{prob:pcp} is undecidable, by reduction from the halting problem.
If there is no solution for given \(\alpha\) and \(\beta\), an algorithm will not halt.

\(\PCP \le \UnInsecure\) can be shown by constructing instances, which add words left or right to an intermediate solution using encryption.
The attacker must find the correct execution order and right amount of instances, such that an entity receives encrypted \([x, x]\) and sends \(\fail\).

A satisfying result from \cref{the:rutu,the:uninsecure} is, that while unbounded sessions are undecidable, a bounded number of sessions is decidable!
Therefore, we can recursively enumerate all insecure protocols, as the problem is semi-decidable.
An algorithm could e.g., search for an attack with sessions bounded to one, then two, three, etc.
If the protocol is insecure, the algorithm \emph{will} terminate, even though, this may take a long time.
If the protocol is secure, the algorithm will not terminate.

This is still pretty unrealistic, as it is computationally expensive, as session count increases instance count and therefore input complexity.
Therefore, we can use approximation algorithms.
One idea is to over-estimate the attacker to reduce the problem space.
Another idea is to model protocol properties in first-order Horn Logic.
This is the underlying thought and the basis of \emph{ProVerif}.

\begin{definition}[Horn Logic]
    We redefine the protocol notation in Horn Logic.
    The Dolev-Yao closure is defined as a one-dimensional relation \(d\) on \(\T\).
    A term \(t \in \T\) is defined to be derivable iff \(d(t)\).
    Therefore we have facts such as \(d(\hat k_C)\), as \(\hat k_C \in I\).

    All rules for the DY closure are written as first-order expressions\footnote{The notation changes to be more familiar to traditional logic.}, e.g.,
    \[\forall t : d(t) \rightarrow d(\hash(t)).\]

    Security is equivalent to \(\neg d(\fail)\).
\end{definition}

The advantage of this notation is, that we can use logical rules between, e.g., literals \(x_1, \dots, x_n, y\), such as \((x_1 \wedge \dots \wedge x_n \rightarrow y) \leftrightarrow (\neg x_1 \vee \dots \vee x_n \vee y)\).

\begin{definition}[Horn Translation of Protocol]
    A protocol is written in Horn Logic by specifying rules for each instance in the original protocol.

    For example, the Needham-Schoeder Protocol can be rewritten in the following way.
    \begin{IEEEeqnarray*}{r"r'c'l}
        A_1: & d(\epsilon) & \rightarrow & d(\enca_{k_C}(A, N_A)) \\
        A_2: & \forall y: d(\enca_{k_A}(N_A, y)) & \rightarrow & d(\enca_{k_C}(y)) \\
        B_1: & \forall x: d(\enca_{k_B}(A, x)) & \rightarrow & d(\enca_{k_A}(x, N_B)) \\
        B_2: & d(\brea, N_B) & \rightarrow & d(\fail)
    \end{IEEEeqnarray*}
    Proving, that the Needham-Schoeder protocol is insecure can be done, be proving \(d(\fail)\) by using the rules and assuming \(d(\epsilon, k_A, k_B, k_C, \hat k_C, \brea, A, B, C)\).
\end{definition}

\begin{definition}[Horn Clauses and Formulas]
    We define Horn clauses to be disjunctions of literals, of which there is at most one positive literal.
    E.g., \(x_2\) or \(x_1 \vee \overline x_2 \vee \overline x_3\).
    The latter is equivalent (\(\equiv\)) to \(x_2 \wedge x_3 \rightarrow x_1\).
    We also have \(\overline x_1 \vee \overline x_2 \equiv \overline{x_1 \wedge x_2} \equiv x_1 \wedge x_2 \rightarrow \false\).

    A Horn formula is a conjunction of Horn clauses.
\end{definition}

The idea is basically to write derivation steps and rules as Horn clauses and evaluate a horn formula.
For this, the lecture additionally defines unifiers, resolvents, and resolutions, and states defines an algorithm to test a horn formula on satisfiability.
Because this is not mentioned any further, and there are no proofs on the aforementioned theorem, this content is not included in this summary.

\begin{center}
    \todo[inline]{Maybe elaborate further.}
\end{center}

\lstset{style=proverif}

\subsection{ProVerif}
ProVerif is a security analysis tool which can be used to prove different securities of protocols, even with unbounded protocol sessions, albeit non-deterministic.

The following are some basic keywords and their semantics.

\begin{lstlisting}
# Free typed algebraic variables
free var
# Communication channel Type
ch: channel
# Data Type
data: bitstring
# Variable the adversary can not directly access.
private
# Security property
query
# Send on channel
out
\end{lstlisting}
